\chapter{Khoảnh khắc vàng}
\markboth{Khoảnh khắc vàng}{Golden Moments}


\section{Em có thể trả lời sau — nhưng em không thể nghĩ thay.}
\begin{truyen}{Em có thể trả lời sau — nhưng em không thể nghĩ thay.}{Em có thể trả lời sau — nhưng em không thể nghĩ th}
\chuhoa{C}{âu thầy nói câu ấy. Cả lớp im lặng — không ai nói, mỗi người đang nghĩ về điều vừa nghe.}
\textit{(The teacher said that. The class fell silent — no one spoke; everyone was thinking about what they had just heard.)}
\end{truyen}

\begin{baihoc}
Im lặng sau câu nói đúng lúc là dấu hiệu của tư duy.
\textit{(Silence after the right sentence is a sign of thinking.)}
\end{baihoc}

\begin{ghinhoanh}
\textbf{Short English to remember:}

Some sentences make the whole class think in silence.

\end{ghinhoanh}

\begin{tuhocnhanh}
1. Câu này khiến em nghĩ đến điều gì? \textit{What does this make you think of?}

2. Em có thể dùng câu này trong tình huống nào? \textit{When could you use this sentence?}
\end{tuhocnhanh}
\ngancach


\section{Im lặng là cách em nói: tôi đang nghĩ.}
\begin{truyen}{Im lặng là cách em nói: tôi đang nghĩ.}{Im lặng là cách em nói: tôi đang nghĩ.}
\chuhoa{C}{âu thầy nói câu ấy. Cả lớp im lặng — không ai nói, mỗi người đang nghĩ về điều vừa nghe.}
\textit{(The teacher said that. The class fell silent — no one spoke; everyone was thinking about what they had just heard.)}
\end{truyen}

\begin{baihoc}
Im lặng sau câu nói đúng lúc là dấu hiệu của tư duy.
\textit{(Silence after the right sentence is a sign of thinking.)}
\end{baihoc}

\begin{ghinhoanh}
\textbf{Short English to remember:}

Some sentences make the whole class think in silence.

\end{ghinhoanh}

\begin{tuhocnhanh}
1. Câu này khiến em nghĩ đến điều gì? \textit{What does this make you think of?}

2. Em có thể dùng câu này trong tình huống nào? \textit{When could you use this sentence?}
\end{tuhocnhanh}
\ngancach


\section{Câu nói khiến cả lớp im là câu nói có sức nặng.}
\begin{truyen}{Câu nói khiến cả lớp im là câu nói có sức nặng.}{Câu nói khiến cả lớp im là câu nói có sức nặng.}
\chuhoa{C}{âu thầy nói câu ấy. Cả lớp im lặng — không ai nói, mỗi người đang nghĩ về điều vừa nghe.}
\textit{(The teacher said that. The class fell silent — no one spoke; everyone was thinking about what they had just heard.)}
\end{truyen}

\begin{baihoc}
Im lặng sau câu nói đúng lúc là dấu hiệu của tư duy.
\textit{(Silence after the right sentence is a sign of thinking.)}
\end{baihoc}

\begin{ghinhoanh}
\textbf{Short English to remember:}

Some sentences make the whole class think in silence.

\end{ghinhoanh}

\begin{tuhocnhanh}
1. Câu này khiến em nghĩ đến điều gì? \textit{What does this make you think of?}

2. Em có thể dùng câu này trong tình huống nào? \textit{When could you use this sentence?}
\end{tuhocnhanh}
\ngancach


\section{Không phải lúc nào nói nhiều cũng là đóng góp — đôi khi im lặng mới là.}
\begin{truyen}{Không phải lúc nào nói nhiều cũng là đóng góp — đôi khi im l...}{Không phải lúc nào nói nhiều cũng là đóng góp — đô}
\chuhoa{C}{âu thầy nói câu ấy. Cả lớp im lặng — không ai nói, mỗi người đang nghĩ về điều vừa nghe.}
\textit{(The teacher said that. The class fell silent — no one spoke; everyone was thinking about what they had just heard.)}
\end{truyen}

\begin{baihoc}
Im lặng sau câu nói đúng lúc là dấu hiệu của tư duy.
\textit{(Silence after the right sentence is a sign of thinking.)}
\end{baihoc}

\begin{ghinhoanh}
\textbf{Short English to remember:}

Some sentences make the whole class think in silence.

\end{ghinhoanh}

\begin{tuhocnhanh}
1. Câu này khiến em nghĩ đến điều gì? \textit{What does this make you think of?}

2. Em có thể dùng câu này trong tình huống nào? \textit{When could you use this sentence?}
\end{tuhocnhanh}
\ngancach


\section{Người học thật là người biết dùng im lặng để suy nghĩ.}
\begin{truyen}{Người học thật là người biết dùng im lặng để suy nghĩ.}{Người học thật là người biết dùng im lặng để suy n}
\chuhoa{C}{âu thầy nói câu ấy. Cả lớp im lặng — không ai nói, mỗi người đang nghĩ về điều vừa nghe.}
\textit{(The teacher said that. The class fell silent — no one spoke; everyone was thinking about what they had just heard.)}
\end{truyen}

\begin{baihoc}
Im lặng sau câu nói đúng lúc là dấu hiệu của tư duy.
\textit{(Silence after the right sentence is a sign of thinking.)}
\end{baihoc}

\begin{ghinhoanh}
\textbf{Short English to remember:}

Some sentences make the whole class think in silence.

\end{ghinhoanh}

\begin{tuhocnhanh}
1. Câu này khiến em nghĩ đến điều gì? \textit{What does this make you think of?}

2. Em có thể dùng câu này trong tình huống nào? \textit{When could you use this sentence?}
\end{tuhocnhanh}
\ngancach


\section{Im lặng cho em cơ hội nghe chính mình.}
\begin{truyen}{Im lặng cho em cơ hội nghe chính mình.}{Im lặng cho em cơ hội nghe chính mình.}
\chuhoa{C}{âu thầy nói câu ấy. Cả lớp im lặng — không ai nói, mỗi người đang nghĩ về điều vừa nghe.}
\textit{(The teacher said that. The class fell silent — no one spoke; everyone was thinking about what they had just heard.)}
\end{truyen}

\begin{baihoc}
Im lặng sau câu nói đúng lúc là dấu hiệu của tư duy.
\textit{(Silence after the right sentence is a sign of thinking.)}
\end{baihoc}

\begin{ghinhoanh}
\textbf{Short English to remember:}

Some sentences make the whole class think in silence.

\end{ghinhoanh}

\begin{tuhocnhanh}
1. Câu này khiến em nghĩ đến điều gì? \textit{What does this make you think of?}

2. Em có thể dùng câu này trong tình huống nào? \textit{When could you use this sentence?}
\end{tuhocnhanh}
\ngancach


\section{Thầy không sợ im lặng — thầy sợ các em nói mà không nghĩ.}
\begin{truyen}{Thầy không sợ im lặng — thầy sợ các em nói mà không nghĩ.}{Thầy không sợ im lặng — thầy sợ các em nói mà khôn}
\chuhoa{C}{âu thầy nói câu ấy. Cả lớp im lặng — không ai nói, mỗi người đang nghĩ về điều vừa nghe.}
\textit{(The teacher said that. The class fell silent — no one spoke; everyone was thinking about what they had just heard.)}
\end{truyen}

\begin{baihoc}
Im lặng sau câu nói đúng lúc là dấu hiệu của tư duy.
\textit{(Silence after the right sentence is a sign of thinking.)}
\end{baihoc}

\begin{ghinhoanh}
\textbf{Short English to remember:}

Some sentences make the whole class think in silence.

\end{ghinhoanh}

\begin{tuhocnhanh}
1. Câu này khiến em nghĩ đến điều gì? \textit{What does this make you think of?}

2. Em có thể dùng câu này trong tình huống nào? \textit{When could you use this sentence?}
\end{tuhocnhanh}
\ngancach


\section{Khoảnh khắc im lặng trong lớp là khoảnh khắc vàng.}
\begin{truyen}{Khoảnh khắc im lặng trong lớp là khoảnh khắc vàng.}{Khoảnh khắc im lặng trong lớp là khoảnh khắc vàng.}
\chuhoa{C}{âu thầy nói câu ấy. Cả lớp im lặng — không ai nói, mỗi người đang nghĩ về điều vừa nghe.}
\textit{(The teacher said that. The class fell silent — no one spoke; everyone was thinking about what they had just heard.)}
\end{truyen}

\begin{baihoc}
Im lặng sau câu nói đúng lúc là dấu hiệu của tư duy.
\textit{(Silence after the right sentence is a sign of thinking.)}
\end{baihoc}

\begin{ghinhoanh}
\textbf{Short English to remember:}

Some sentences make the whole class think in silence.

\end{ghinhoanh}

\begin{tuhocnhanh}
1. Câu này khiến em nghĩ đến điều gì? \textit{What does this make you think of?}

2. Em có thể dùng câu này trong tình huống nào? \textit{When could you use this sentence?}
\end{tuhocnhanh}
\ngancach


\section{Suy nghĩ không ồn ào — nhưng nó thay đổi con người.}
\begin{truyen}{Suy nghĩ không ồn ào — nhưng nó thay đổi con người.}{Suy nghĩ không ồn ào — nhưng nó thay đổi con người}
\chuhoa{C}{âu thầy nói câu ấy. Cả lớp im lặng — không ai nói, mỗi người đang nghĩ về điều vừa nghe.}
\textit{(The teacher said that. The class fell silent — no one spoke; everyone was thinking about what they had just heard.)}
\end{truyen}

\begin{baihoc}
Im lặng sau câu nói đúng lúc là dấu hiệu của tư duy.
\textit{(Silence after the right sentence is a sign of thinking.)}
\end{baihoc}

\begin{ghinhoanh}
\textbf{Short English to remember:}

Some sentences make the whole class think in silence.

\end{ghinhoanh}

\begin{tuhocnhanh}
1. Câu này khiến em nghĩ đến điều gì? \textit{What does this make you think of?}

2. Em có thể dùng câu này trong tình huống nào? \textit{When could you use this sentence?}
\end{tuhocnhanh}
\ngancach


\section{Em có quyền không trả lời ngay — và dùng quyền đó để nghĩ.}
\begin{truyen}{Em có quyền không trả lời ngay — và dùng quyền đó để nghĩ.}{Em có quyền không trả lời ngay — và dùng quyền đó }
\chuhoa{C}{âu thầy nói câu ấy. Cả lớp im lặng — không ai nói, mỗi người đang nghĩ về điều vừa nghe.}
\textit{(The teacher said that. The class fell silent — no one spoke; everyone was thinking about what they had just heard.)}
\end{truyen}

\begin{baihoc}
Im lặng sau câu nói đúng lúc là dấu hiệu của tư duy.
\textit{(Silence after the right sentence is a sign of thinking.)}
\end{baihoc}

\begin{ghinhoanh}
\textbf{Short English to remember:}

Some sentences make the whole class think in silence.

\end{ghinhoanh}

\begin{tuhocnhanh}
1. Câu này khiến em nghĩ đến điều gì? \textit{What does this make you think of?}

2. Em có thể dùng câu này trong tình huống nào? \textit{When could you use this sentence?}
\end{tuhocnhanh}
\ngancach
